\section{题02：圆形灯环（P16784）}

\subsection{问题建模：灯环约束转化为有向图上的闭游走}

一个由$N=2026$个LED灯珠首尾相连组成的圆形灯环，每个灯珠的亮度可在$0\sim 25$共$K=26$个等级中独立选择。一个配置“稳定”需满足两个条件：顺时针方向上任意连续$3$个灯珠的亮度之和$\neq 26$；相邻对$(2,6)$（前一个为$2$级、后一个为$6$级）在整个环中出现总次数为偶数。能通过顺时针旋转相互重合的配置视为本质相同。要求计算本质不同的稳定配置数模$998244353$的结果。

核心思路是将环形序列建模为有向图上的闭游走。定义状态为相邻两灯珠的亮度对$(x,y)$，共有$K^2=26^2=676$个状态，每个状态用编号$x\cdot K+y$唯一标识。状态$(x,y)$可一步转移到$(y,z)$当且仅当三元组约束不触犯，即$x+y+z\neq 26$。例如状态$(5,10)$：枚举$z=0\sim 25$，$5+10+z=26$解得$z=11$，故$z=11$被禁止，其余$25$个$z$值均可转移；状态$(0,0)$：$0+0+z=26$要求$z=26$超出范围，因此全部$26$个$z$均可转移；状态$(12,14)$：$12+14+z=26$给出$z=0$，仅$z=0$被禁止。每个状态的出度在$25$或$26$之间。

考虑一个长度为$d$的环形配置$a_0,a_1,\dots,a_{d-1}$（下标模$d$）。以$(a_0,a_1)$为起点，依次转移至$(a_1,a_2),\dots,(a_{d-1},a_0)$，最终回到$(a_0,a_1)$，这恰好是一条长度为$d$的闭游走。游走的$d$次转移依次检验了环上全部$d$个三元组约束，且访问的$d$个状态正是环上$d$个相邻对。关键观察：从形序列到闭游走的映射是一一且映上的——给定闭游走可唯一还原环形序列$a_i$，给定环形序列可唯一构造闭游走。因此数环形配置等价于数闭游走。

\subsection{朴素思路：$26^{2026}$的枚举与为什么需要矩阵}

最直接的暴力算法是枚举全部$26^{2026}$种亮度分配，逐一验证三元组约束与$(2,6)$奇偶性，再用循环移位去重。$26^{2026}\approx 10^{2866}$远超任何计算机能处理的规模——整个可观测宇宙的原子数也仅约$10^{80}$。即便采用回溯法逐位搜索并利用三元组约束在每一位剪枝，$2026$位的搜索深度仍完全无法承受。然而约束的局部性——三元组只涉及连续三个灯珠，$(2,6)$计数可沿途累积——暗示状态空间可压缩为基于“最近两个灯珠”的马尔可夫链，进而用矩阵乘法批量计算所有路径。这正是从指数级暴力到多项式级代数计算的关键跳板。

\subsection{核心洞察：邻接矩阵的迹即环形序列计数}

将图中允许的转移记为$676\times 676$邻接矩阵$M$：$M_{ij}=1$若状态$i$可一步转移到状态$j$，否则为$0$。由矩阵乘法的组合解释，$(M\times M)_{ij}=\sum_k M_{ik}M_{kj}$枚举了从$i$经中间状态$k$到$j$的所有两步路径，恰好对应两步游走的组合计数。归纳地，$M^d$的第$(i,j)$元等于从$i$到$j$的长度为$d$的游走总数。

因此$\mathrm{tr}(M^d)=\sum_i (M^d)_{ii}$给出了所有长度为$d$的闭游走总数，记为$A(d)$。以一个$d=3$的具体环验证：设三个灯珠亮度为$(2,5,8)$，闭游走为$(2,5)\to(5,8)\to(8,2)\to(2,5)$。第一步检验$2+5+8=15\neq 26$，第二步检验$5+8+2=15\neq 26$，第三步检验$8+2+5=15\neq 26$。三步恰好覆盖环上全部$3$组连续三元组。对于$d=1$的平凡情形，$A(1)=\mathrm{tr}(M)$仅统计所有形如$(x,x)$的自环，要求$x+x+x\neq 26$。$3x=26$无整数解，故全部$26$个自环均合法，$A(1)=26$。

\subsection{符号矩阵：以$-1$权值从代数上分离$(2,6)$奇偶性}

$(2,6)$出现次数的奇偶性约束通过一个精巧的代数技巧解决。构造带符号矩阵$M'$：若某条转移的目标状态恰好为$(2,6)$（即$y=2,z=6$），则边权取$-1$（模意义下置为$\mathrm{MOD}-1$），否则取$1$。考虑任一长度为$d$的闭游走：目标为$(2,6)$的转移每出现一次贡献一个$-1$因子，其余转移贡献$1$，因此沿游走的全部边权乘积精确等于$(-1)^{\#(2,6)}$。这正是生成函数的思想——用$-1$给$(2,6)$的出现加权，在求和层面实现了奇偶性的差分编码。

故$B(d)=\mathrm{tr}((M')^d)$等于偶数情形数$C_0(d)$减去奇数情形数$C_1(d)$，即$B(d)=C_0(d)-C_1(d)$。同时$A(d)=\mathrm{tr}(M^d)=C_0(d)+C_1(d)$给出了不考虑奇偶性的所有合法环形序列数。两式相加得$2C_0(d)=A(d)+B(d)$，故$C_0(d)=(A(d)+B(d))/2$。模$998244353$下的除法乘以$2$的模逆元$(998244353+1)/2=499122177$实现。

\subsection{周期约束归约：何时用$A(d)$、何时用$C_0(d)$}

现在将视线转到Burnside引理下的$h(d)$定义上。$h(d)$是周期整除$d$的、长度为$N$的合法配置总数。周期为$d$的一个$N$长配置由一段长为$d$的“基本模式”重复$r=N/d$次铺排而成。在此重复构造下，环上全部三元组约束等价于要求该基本模式自身构成一个长度为$d$的合法环形序列——因为所有$d$个本质不同的三元组在重复次数增加时并不产生新的约束形式，且首尾衔接处也由基本模式的闭合性自动保证。

$(2,6)$的出现情况则需分情况讨论。设基本模式中$(2,6)$出现$t$次（$0\le t\le d$），则完整环上共出现$r\times t$次。若$r$为偶数，$r\times t$恒为偶数，与$t$无关，此时$h(d)$直接取$A(d)$；若$r$为奇数，则要求$t$本身为偶数，故需用$C_0(d)$。这一推理的优雅之处在于，它将一个全局约束（$N=2026$长度上的奇偶性）约化为了基本模式长度$d$上的局部约束，从而可以与矩阵幂的结果无缝对接。

\subsection{Burnside求和：四项因子的具体展开}

$N=2026=2\times 1013$，正因子集为$\{1,2,1013,2026\}$。四项的重复次数$r=N/d$依次为$2026,1013,2,1$。

由前述规则判定每一项的$h(d)$。$r=2026$为偶数$\Rightarrow h(1)=A(1)=26$。$r=1013$为奇数$\Rightarrow h(2)=C_0(2)$。$r=2$为偶数$\Rightarrow h(1013)=A(1013)$。$r=1$为奇数$\Rightarrow h(2026)=C_0(2026)$。

欧拉函数值的计算：$\varphi(2026)=\varphi(2)\varphi(1013)=1\times 1012=1012$（$1013$为素数故$\varphi(1013)=1012$）；$\varphi(1013)=1012$；$\varphi(2)=1$；$\varphi(1)=1$。

代入Burnside公式，所有项除以$N=2026$并取模：
\[
\mathrm{ans}=\frac{1012\cdot 26+1012\cdot C_0(2)+1\cdot A(1013)+1\cdot C_0(2026)}{2026}\pmod{998244353}.
\]
式中的除法通过乘以$2026$在模$998244353$下的逆元完成。由于$998244353$是素数且$2026$与之互素，该逆元必然存在。

\subsection{矩阵快速幂：二进制分解与访存优化}

为计算$A(1013)=\mathrm{tr}(M^{1013})$和$C_0(2026)$等项，需高效计算高次矩阵幂。直接迭代$1013$次矩阵乘法，单次$676^3\approx 3.09\times 10^8$次操作，总计$\sim 3\times 10^{11}$，不可接受。采用二进制快速幂策略：预处理$M^{2^k}$（$k=0,\dots,9$即$M^1$至$M^{512}$）共$9$次矩阵乘法。$1013$的二进制为$1111110101_2=512+256+128+64+32+16+4+1$，从$M^1$出发依次乘上$M^{2^2}$至$M^{2^9}$的$7$个已置位方幂，共$7$次乘法得到$M^{1013}$。

$\mathrm{tr}(M^{2026})$利用恒等式$\mathrm{tr}(M^{2026})=\sum_{i,j}(M^{1013})_{ij}(M^{1013})_{ji}$直接计算——对$M^{1013}$做$i,j$的二重循环累加$M_{ij}\cdot M_{ji}$即可，无需完整矩阵乘法。对$M'$同理，但$2026=1024+512+256+128+64+32+8+2$需预处理至$M'^{1024}$（$k=0,\dots,10$共$11$个方幂、$10$次乘法），再从$M'^2$出发乘上$7$个方幂得到$M'^{2026}$。

稠密矩阵乘法的性能瓶颈在于访存。为提升缓存局部性，预先转置右乘矩阵$B$为$B^T$。$C_{ij}=\sum_k A_{ik}B^T_{jk}$在内层对$k$的循环中，$A_{ik}$和$B^T_{jk}$均为行主序连续访问，大幅减少cache miss。累加过程使用\mil{__int128}类型容纳中间和——最坏情况约$676\times (998244352)^2\approx 6.76\times 10^{20}$，远低于\mil{__int128}上限$1.7\times 10^{38}$，遍历完一行$k$后对每个$C_{ij}$取模一次。对$M$和$M'$各自约需$17$次稠密矩阵乘法，总计约$34$次，每次$3.09\times 10^8$次基本操作，总操作量约$1.05\times 10^{10}$量级，在现代CPU上数秒内完成。空间需存储约$10$个$676\times 676$矩阵及对应转置，约$10\times 676^2\times 4\times 2\approx 36$MB，在合理范围内。

\subsection{代码实现}

本题为填空题，实际提交只需输出预计算的答案。但为了完整展示求解过程，以下同时给出幕后求解代码与提交代码。

\medskip\noindent 求解代码：
\inputminted{c++}{../src/p02_circular_light_ring_solve.cpp}

\medskip\noindent 提交代码：
\inputminted{c++}{../src/p02_circular_light_ring.cpp}
